\section{The synthetic correspondence instrument}
\label{sec:synthetic-instrument}

No measurement of correspondence, as distinct from surface coverage, existed
in the pipeline before this internship. One was built from synthetic data
generated by the model itself: parameters $(\bm\beta, \bm\theta)$ are sampled
from the model's own learned distributions, the template is posed
accordingly, the result is exported as an anonymous mesh, optionally
degraded with noise, holes, or partial occlusion to resemble real scan
artefacts, and passed back through the fitter exactly as a real specimen
would be. Because the mesh originates from the model, the true identity of
every vertex is known by construction, so per-vertex correspondence error is
directly measurable rather than inferred indirectly from surface distance.
This instrument underlies every diagnostic in the remainder of this chapter.

\begin{investigation}{I1}{Is correspondence error a part-identity problem or a localisation problem?}
\invQuestion{When a vertex lands in the wrong place, would telling the
optimiser which anatomical region it belongs to remove most of the error?}
\invMethod{Round-trip fits of synthetic specimens with known vertex identity,
comparing the default fit against a fit given the correct \emph{regional
identity} of every target point (a counterfactual, or \emph{oracle},
intervention; \S\ref{sec:oracle}).}
\invResult{Leg accuracy improves from $0.867$ to $0.937$: a real but partial
recovery. Substantial error remains \emph{inside} correctly identified
structures.}
\invVerdict{DIAGNOSTIC}{Regional identity and within-region localisation are
separate problems, and fixes aimed only at part discrimination (segmentation
refinements, SDF-style descriptors) address only the first. This is a
causal-bottleneck finding, not a shipped intervention: no algorithm here is
adopted into the pipeline. Recorded as \F{F3}. An earlier percentage for this
split was derived under a reference later found to be invalid (\F{F7}) and is
not reported.}
\end{investigation}
